\section{ERD hệ thống OLTP và luồng dữ liệu}\label{sec:erd-oltp}
Sơ đồ quan hệ thực thể của hệ thống bán hàng có \code{orders} làm trung tâm: mỗi khách (\code{customers}) đặt nhiều đơn; mỗi đơn có nhiều dòng hàng (\code{order\_\allowbreak{}items}), nhiều lần thanh toán (\code{order\_\allowbreak{}payments}) và nhiều đánh giá (\code{order\_\allowbreak{}reviews}). Dòng hàng nối tới sản phẩm (\code{products}) và người bán (\code{sellers}). Toàn bộ nguồn có 11 khoá chính và 7 khoá ngoại. Khoá chính của \code{order\_\allowbreak{}items} là cặp (\code{order\_\allowbreak{}id}, \code{order\_\allowbreak{}item\_\allowbreak{}id}), của \code{order\_\allowbreak{}payments} là (\code{order\_\allowbreak{}id}, \code{payment\_\allowbreak{}sequential}). Với \code{order\_\allowbreak{}reviews}, mã \code{review\_\allowbreak{}id} xuất hiện lặp ở nhiều đơn nên không dùng làm khoá chính được; khoá là cặp (\code{review\_\allowbreak{}id}, \code{order\_\allowbreak{}id}). Bảng \code{geolocation} không có khoá tự nhiên, vì một mã bưu chính có nhiều toạ độ, nên được thêm cột kỹ thuật \code{id}; nó chỉ liên hệ logic với khách và người bán qua mã bưu chính. Cũng chỉ liên hệ logic là \code{products} và bảng dịch tên danh mục, vì có danh mục không có bản dịch. Hệ thống CRM gồm \code{mql} và \code{closed\_\allowbreak{}deals} nối theo quan hệ 1–0..1 (một khách tiềm năng có thể được chốt thành một thương vụ). Cột \code{closed\_\allowbreak{}deals.seller\_\allowbreak{}id} nối logic sang \code{oltp.sellers} bằng đường nét đứt trong sơ đồ và không có khoá ngoại. Mã bưu chính trong bản dữ liệu này đã đầy đủ 5 chữ số, kể cả số 0 đầu, nên được lưu dạng văn bản.

\hinh{3-2-4_erd_oltp.png}{ERD hai hệ thống nguồn: bán hàng (oltp) và marketing (crm)}{fig:erd-oltp}{1.0}
\todo{Yêu cầu của môn: trên ERD OLTP \textbf{gạch chân các trường dữ liệu sẽ dùng sang OLAP} và ghi chú \textbf{mapping} tương ứng. Vẽ lại ERD bằng draw.io từ hình nháp, gạch chân các cột dùng và ghi chú (bảng mapping đầy đủ ở Phụ lục C, nguồn docs/mapping\_oltp\_olap.md). Thêm sơ đồ luồng dữ liệu (data flow) từ nguồn đến kho.}

\subsection{ERD tô màu và các cột dùng sang OLAP}\label{sec:erd-mau}
Hình \ref{fig:erd-mau} vẽ lại ERD với màu theo nhóm dữ liệu (dữ liệu chủ, giao dịch, tương tác, tham chiếu, cùng nhóm với cây phân loại ở mục \ref{sec:taxonomy}); các cột \underline{gạch chân} là cột nguồn được dùng để tạo cột của dimension hoặc fact, suy ra tự động từ bảng mapping (Phụ lục C). Trong 71 cột của 11 bảng nguồn có 65 cột được dùng sang OLAP; 6 cột còn lại gồm khoá kỹ thuật \code{geolocation.id}, ba thuộc tính mô tả sản phẩm không dùng (độ dài tên, độ dài mô tả, số ảnh) và hai thuộc tính khai báo của deal (doanh thu hằng tháng và quy mô danh mục sản phẩm khai báo). Bảy thuộc tính còn lại của deal được gom vào chiều rác \code{dim\_lead\_profile}.
\hinh{2-5_erd_oltp_gach_chan.png}{ERD hệ thống nguồn tô màu theo nhóm dữ liệu, gạch chân các cột dùng sang OLAP (bản nháp)}{fig:erd-mau}{1.0}

\subsection{Tổng hợp mapping OLTP sang OLAP}\label{sec:mapping-tong-hop}
Bảng \ref{tab:mapping-tong-hop} tổng hợp, cho từng bảng nguồn, số cột được dùng, các bảng dimension và fact nhận dữ liệu và một ví dụ phép biến đổi; bảng mapping đầy đủ từng cột ở Phụ lục C. Hình \ref{fig:flow-oltp-olap} cho thấy luồng dữ liệu từ bảng nguồn qua ODS đến từng dimension và fact.
\begin{table}[htbp]\centering\scriptsize\caption{Tổng hợp mapping OLTP sang OLAP theo bảng nguồn}\label{tab:mapping-tong-hop}
\begin{tabular}{p{2.68cm}p{0.89cm}p{1.25cm}p{4.47cm}p{4.29cm}}\toprule
\textbf{Bảng nguồn} & \textbf{Số cột} & \textbf{Cột dùng sang OLAP} & \textbf{Bảng đích (dim, fact)} & \textbf{Ví dụ phép biến đổi}\\\midrule
oltp.customers & 5 & 5 & dim\_customer, dim\_geography, fact\_order\_fulfillment, fact\_order\_item, fact\_payment, fact\_review & customer\_zip\_code\_prefix $\rightarrow$ zip\_code\_prefix: Chuẩn hoá char(5): bỏ ký tự không phải số, lpad(…,5,'0')\\
oltp.geolocation & 6 & 5 & dim\_geography & geolocation\_zip\_code\_prefix $\rightarrow$ zip\_code\_prefix: Chuẩn hoá char(5) rồi GROUP BY zip\\
oltp.order\_items & 7 & 7 & fact\_order\_fulfillment, fact\_order\_item & giữ nguyên giá trị\\
oltp.order\_payments & 5 & 5 & fact\_order\_fulfillment, fact\_payment & payment\_type $\rightarrow$ payment\_type: lower(trim())\\
oltp.order\_reviews & 7 & 7 & fact\_review & review\_id $\rightarrow$ review\_id: Khử trùng trên cặp (review\_id, order\_id): giữ bản có review\_answer\_timestamp mới nhất, dòng còn lại ghi etl\_log.rejected\_rows\\
oltp.orders & 8 & 8 & dim\_customer, fact\_order\_fulfillment, fact\_order\_item, fact\_payment, fact\_review & order\_status $\rightarrow$ order\_status: lower(trim())\\
oltp.product\_category\_name\_translation & 2 & 2 & dim\_product & giữ nguyên giá trị\\
oltp.products & 9 & 6 & dim\_product & product\_category\_name $\rightarrow$ product\_category\_name: trim; rỗng/NULL $\rightarrow$ 'unknown'\\
oltp.sellers & 4 & 4 & dim\_geography, dim\_seller & seller\_zip\_code\_prefix $\rightarrow$ zip\_code\_prefix: Chuẩn hoá char(5): bỏ ký tự không phải số, lpad(…,5,'0')\\
crm.closed\_deals & 14 & 8 & dim\_lead\_profile, dim\_sales\_rep, dim\_seller, fact\_lead\_funnel & business\_segment $\rightarrow$ business\_segment: trim; NULL/rỗng $\rightarrow$ 'unknown'\\
crm.mql & 4 & 4 & dim\_seller, fact\_lead\_funnel & origin $\rightarrow$ origin: lower(trim()); NULL/rỗng $\rightarrow$ 'unknown'\\
\midrule
\textbf{Tổng} & 71 & 61 & & \\
\bottomrule\end{tabular}\end{table}

\hinh{2-5_data_flow_oltp_olap.png}{Luồng dữ liệu OLTP sang OLAP: bảng nguồn nào nuôi bảng nào của kho (bản nháp)}{fig:flow-oltp-olap}{1.0}
\todo{Nhóm vẽ lại ERD bằng draw.io từ docs/diagrams/erd-oltp-tomau.drawio (giữ tên file hình 2-5\_erd\_oltp\_gach\_chan.png) và sơ đồ data flow OLTP sang OLAP từ docs/diagrams/data-flow-oltp-olap.drawio.}
